\documentclass[12pt,letterpaper]{article}
\usepackage[utf8]{inputenc}
\usepackage{amsmath}
\usepackage{amsfonts}
\usepackage{amssymb}
\usepackage{amsthm}
\usepackage{graphicx}
\usepackage{tabularx}
\usepackage[left=2cm,right=2cm,top=2cm,bottom=2cm]{geometry}
\usepackage{fancyhdr}
\usepackage{multicol}
\usepackage{multirow,array}
\usepackage{newtxtext,newtxmath}
\usepackage{relsize}
\usepackage{lastpage}
\usepackage{cancel}
\usepackage{tikz}
\usepackage{enumitem}
\usepackage{adjustbox}
\newcolumntype{Y}{>{\centering\arraybackslash}X}
	\setenumerate[1]{label={\bf \theenumi: ~}}
	\setenumerate[2]{label={\bf \theenumii: ~}}
\pagestyle{fancy}
\fancyhf{}
\lhead{\textsc{BHCC Mat-181}}
\rhead{\textsc{Normal areas and $z$ scores}}
\rfoot{Page \thepage ~of \pageref{LastPage}}

\newcommand*\circled[1]{\tikz[baseline=(char.base)]{
            \node[shape=circle,draw,inner sep=2pt] (char) {#1};}}
\newcommand{\N}[2]{\mathcal{N}\big(#1,~#2\big)}

\begin{document}
\begin{enumerate}
\item Let $X\sim \N{400}{10}$. You will calculate $P(414 < X < 422)$ after completing the diagram.
\begin{center}
\includegraphics[scale=.8]{figures/curve.png}
\end{center}
\begin{enumerate}
\item Determine the $x$ values that correspond to the integer $z$ values.
\item Find the $z$ scores of the boundaries, and shade the appropriate section.
\item Estimate $P(414 < X < 422)$ by counting squares.
\vspace{30pt}
\item Calculate $P(414 < X < 422)$ by using a $z$ table.
\end{enumerate}
 \vfill
 
\item Let $X\sim \N{5.5}{0.1}$. You will calculate $P\Big(\big|X-5.5 \big| > 0.13 \Big)$ after completing the diagram.
\begin{center}
\includegraphics[scale=.8]{figures/curve.png}
\end{center}
\begin{enumerate}
\item Determine the $x$ values that correspond to the integer $z$ values.
\item Find the $z$ scores of the boundaries, and shade the appropriate section.
\item Estimate $P\Big(\big|X-5.5 \big| > 0.13 \Big)$ by counting squares.
\vspace{30pt}
\item Calculate $P\Big(\big|X-5.5 \big| > 0.13 \Big)$ by using a $z$ table.
\end{enumerate}
\vfill

\newpage

\item Let $X \sim \N{23.4}{5.6}$. Determine $P(X>21)$.
\vfill
\item Let $X \sim \N{100}{5}$. Determine $x_0$ such that $P(X<x_0) = 0.30$.
\vfill
\item Let $X \sim \N{100}{5}$. Determine $x_1$ such that $P(X>x_1) = 0.30$.
\vfill
\item Let $X \sim \N{100}{5}$. Determine $r$ such that $P\big(|X-100|<r\big) = 0.30$.
\vfill
\item Let $X \sim \N{100}{5}$. Determine $r$ such that $P\big(|X-100|>r\big) = 0.30$.
\vfill

\newpage

\item Determine $\sigma$ such that $X \sim \N{10}{\sigma}$ and $P(X > 12.3) = 0.42$.
\vfill
\item Let $X \sim \N{42.5}{0.15}$. Determine $P(|X-42.5| < 0.25)$.
\vfill
\item Determine $\mu$ such that $X \sim \N{\mu}{30}$ and $P(X > 150) = 0.8$.
\vfill
\item Let $X \sim \N{1000}{30}$. Determine $P(930 < X < 991)$.
\vfill
\end{enumerate}

\newpage
\chead{\bf ANSWERS}
\begin{multicols}{2}
\begin{enumerate}
\item \begin{enumerate} \item 370, 380, 390, 400, 410, 420, 430
\item $z_1 = \frac{414-400}{10} = 1.4$\\
$z_2 = \frac{422-400}{10} = 2.2 $\\
\includegraphics[scale=0.4]{figures/A1.png}
\item $P(414<X<422) \approx 7\%$
\item $\Phi(2.2)-\Phi(1.4) = \fbox{0.06685}$
\end{enumerate}

\item \begin{enumerate} \item 5.2, 5.3, 5.4, 5.5, 5.6, 5.7, 5.8
\item $z_1 = \frac{5.37-5.5}{0.1} = -1.3$\\
$z_2 = \frac{5.63-5.5}{0.1} = 1.3 $\\
\includegraphics[scale=0.4]{figures/A2.png}
\item $P(|X-5.5|<0.13) \approx 18\%$\\
$P(X<5.37 \textsc{~~or~~} X>5.63) \approx 18\%$

\item $2\Phi(-1.3) = \fbox{0.1936}$
\end{enumerate}

\item $z = \frac{21-23.4}{5.6} = -0.43$\\
$P(X>21) = P(Z>-0.43)=$\\ 
$= 1-\Phi(-0.43) ~=~ \fbox{0.666}$

\item We recognize $x_0$ is 30th percentile. \\
$z = \Phi^{-1}(0.3) = -0.52$\\
$x_0 = z\sigma+\mu = (-0.52)(5)+100 = \fbox{97.4}$

\item We recognize $x_1$ is 70th percentile. \\
$z = \Phi^{-1}(0.7) = 0.52$\\
$x_0 = z\sigma+\mu = (0.52)(5)+100 = \fbox{102.6}$

\item This is harder. We recognize $(100+r)$ is the 65th percentile. We get 65 by splitting 30 in half and adding it to 50. A sketch will help here. \\
$z_{\textsc{upper}} = \Phi^{-1}(0.65) = 0.38$\\
$x_{\textsc{upper}} = (0.38)(5)+100 = 101.9$\\
$r = \fbox{1.9}$

\item We recognize $(100+r)$ is the 85th percentile. We get 85 by splitting 30 in half and subtracting it from 100. A sketch will help here. \\
$z_{\textsc{upper}} = \Phi^{-1}(0.85) = 1.04$\\
$x_{\textsc{upper}} = (1.04)(5)+100 = 105.2$\\
$r = \fbox{5.2}$

\item 42\% is a right area corresponding to 58th percentile. We find the $z$ score of 58th percentile.\\
$z = \Phi^{-1}(0.58) = 0.20 $\\
$\sigma = \frac{x-\mu}{z} =\frac{12.3-10}{0.2} = \fbox{11.5}$

\item We want a central area with bounds 42.25 and 42.75.\\
$z_{\textsc{lower}} = \frac{-0.25}{0.15} = -1.67$ \\
$z_{\textsc{upper}} = \frac{0.25}{0.15} = 1.67$ \\
$P(|X-42.5|<0.25) = \Phi(1.67)-\Phi(-1.67) = \fbox{0.905}$

\item We want a $z$ score of 20th percentile.\\
$z = \Phi^{-1}(0.2) = -0.84$ \\
$\mu = x - z\sigma = 150-(-0.84)(30) = \fbox{175.2}$

\item We find two $z$ scores.\\
$z_{\textsc{lower}} = \frac{930-1000}{30} = -2.33$\\
$z_{\textsc{upper}} = \frac{991-1000}{30} = -0.3$ \\
$P(930<X<991) = \Phi(-0.3)-\Phi(-2.33) = \fbox{0.3722}$

\end{enumerate}
\end{multicols}
\end{document}

